\originalchapter{作业 4}
MIT 18.330, Spring 2012. 原件: Homework 4, Due Tuesday April 3. 以下解答均为 AI 编写, 非官方答案.

\begin{exercise}\label{ps4-q1}
(原题 1, 1 分) 用 Newton 方法求 $3^{1/3}$, 达到 10 位精度, 并说明计算过程.
\end{exercise}
\begin{solution}
解 $f(x)=x^3-3=0$, 取 $x_0=1$, 迭代
$x_{k+1}=(2x_k+3/x_k^2)/3$. 前几项为
$1.6666666667$, $1.4711111111$, $1.4428120982$, $1.4422497896$,
$1.4422495703074416$. 下一步得到 $1.4422495703074085$, 残差约 $8.9\times10^{-16}$. 因此小数点后 10 位为 $1.4422495703$. 用 $f(1.4422495703)<0<f(1.44224957035)$ 及 $f'>0$ 可复核根的区间; 脚本保存实际迭代值.
\end{solution}

\begin{exercise}\label{ps4-q2}
(原题 2) 解 $x=\phi(x)$ 的一种方法是迭代 $x_{k+1}=\phi(x_k)$.
\begin{enumerate}[label=(\alph*)]
\item (1 分) 当 $\phi$ 为压缩映射, 即对所有 $x\ne y$ 有
$|\phi(x)-\phi(y)|<|x-y|$, 迭代收敛. 证明若对所有 $x$ 有 $|\phi'(x)|<1$, 则满足这个压缩性质.
\item (0.5 分) 找一个 $\phi$, 使固定点唯一, 但迭代发散.
\item (1 分) 设 $f$ 只有一个根 $x_*$, 且在根的邻域内 $f'\ne0$. 将 Newton 法写成固定点迭代, 用 (a) 给出一个关于 $f,f',f''$ 的邻域条件, 保证收敛到固定点.
\end{enumerate}
\end{exercise}
\begin{solution}
(a) 在定义域为区间且 $\phi$ 可微时, 中值定理给出
$|\phi(x)-\phi(y)|=|\phi'(\xi)||x-y|<|x-y|$.
整理注: 原题的逐对严格不等式比通常 Banach 压缩条件弱. Banach 定理的一种充分保证是不变的完整区间及统一常数 $q<1$, 不能仅凭逐点 $|\phi'|<1$ 断言任意起点收敛. 例如 $\phi(x)=\log(1+\ee^x)$ 的导数处处在 $(0,1)$ 内, 但没有实固定点.

(b) 取 $\phi(x)=2x$, 唯一固定点为 $0$; 任意 $x_0\ne0$ 的迭代为 $2^kx_0$, 故发散.

(c) $\phi(x)=x-f(x)/f'(x)$,
$\phi'(x)=f(x)f''(x)/[f'(x)]^2$. 充分条件是在闭区间
$I=[x_*-r,x_*+r]$ 上有 $f'\ne0$, $f\in C^2$, 且
$\sup_I|ff''/(f')^2|\le q<1$. 因 $\phi(x_*)=x_*$, 中值定理给出
$|\phi(x)-x_*|\le q|x-x_*|\le qr$, 所以 $\phi(I)\subset I$.
对 $x_0\in I$, 误差不超过 $q^k|x_0-x_*|$, 故迭代收敛.
\end{solution}

\begin{exercise}\label{ps4-q3}
(原题 3, 2.5 分) 用多变量 Newton 法求下列方程组的一个解:
\[
x_1^2+x_2^2+x_3^2=100,\qquad x_1x_2x_3=1,\qquad
x_1-x_2-\sin x_3=0.
\]
\end{exercise}
\begin{solution}
将三个左端减去右端组成 $G(x)$, Jacobian 为
\[
J(x)=\begin{pmatrix}
2x_1&2x_2&2x_3\\
x_2x_3&x_1x_3&x_1x_2\\
1&-1&-\cos x_3
\end{pmatrix}.
\]
每步解 $J(x_k)\Delta x=-G(x_k)$, 再取 $x_{k+1}=x_k+\Delta x$, 不显式求逆. 从 $(7,7,0.02)$ 出发实际得到
\[
(x_1,x_2,x_3)\approx
(7.081046012560867,\ 7.061047185906358,\ 0.02000015999317334).
\]
三方程残差最大绝对值小于 $2\times10^{-14}$. 题目只要求一个解; 该实验不主张方程组解唯一. 源码中的 \texttt{newton} 函数 保留各步残差.
\end{solution}

\begin{exercise}\label{ps4-q4}
(原题 4) 求 $F$ 极小值的 Newton 迭代为
$x_{k+1}=x_k-F'(x_k)/F''(x_k)$. 以下取
$F(x)=1+\int_0^x\arctan y\dd y$.
\begin{enumerate}[label=(\alph*)]
\item (0.5 分) 证明 $F''(x)>0$, 即 $F$ 严格凸. (原件附注: 严格凸函数总有唯一极小值.)
\item (0.5 分) 分别给出一个收敛的起点及一个发散的起点; 凸性不能保证 Newton 迭代收敛.
\item (1 分) 简述如何设计一个可靠方法, 在满足 $F'(a)<0$, $F'(b)>0$ 的区间 $[a,b]$ 中求凸函数的极小值.
\end{enumerate}
\end{exercise}
\begin{solution}
(a) $F'(x)=\arctan x$, $F''(x)=1/(1+x^2)>0$. 本题唯一极小点为 $0$, 极小值为 $1$. 整理注: 一般严格凸性只保证极小点至多一个, 不保证其存在, 如 $\ee^x$ 在 $\R$ 上没有最小点; 本题存在性由 $F'(0)=0$ 及导数变号给出.

(b) 迭代函数为 $\phi(x)=x-(1+x^2)\arctan x$. 取 $x_0=1/2$, 在 $|x|\le1/2$ 上
$|\phi'(x)|=2|x\arctan x|\le\arctan(1/2)<1$, 且 $\phi(0)=0$, 所以收敛.
取 $x_0=2$: 若 $x\ge2$, 令
$H(x)=(1+x^2)\arctan x-2x$; 有 $H(2)>0$,
$H'(x)=2x\arctan x-1>0$, 故 $\phi(x)<-x$.
$\phi$ 为奇函数, 所以迭代符号交替且绝对值严格增加. 若其绝对值有有限极限 $L\ge2$, 连续性将给出
$(1+L^2)\arctan L-L=L$, 与 $H(L)>0$ 矛盾; 故绝对值趋无穷.

(c) 对连续导数 $F'$ 二分. 在中点 $m$ 检查导数符号, 保留含零点的半区间; 区间长度每步减半. 凸性使 $F'$ 单调不减, 零点就是极小点; 若严格凸则极小点唯一. 也可加入受区间约束的 Newton 步, 但必须保留二分的收缩保证.
\end{solution}

\begin{exercise}\label{ps4-q5}
(原题 5, 2 分) 要精确测量电阻, 其标称值为 $R=2\Omega$. 接电池后用廉价万用表测得电压 $V=2.9\mathrm V$, 电流 $I=1.4\mathrm A$. 假定标称值和测量均有误差, 用 Newton 法最小化
\[
F(V,I,R)=(V-RI)^2+10(R-2)^2+10(V-2.9)^2+10(I-1.4)^2,
\]
求 $V,I,R$ 的“最佳”拟合.
\end{exercise}
\begin{solution}
令 $q=(V,I,R)^T$, $q_0=(2.9,1.4,2)^T$, $d=V-RI$,
$g=(1,-R,-I)^T$. 则
\[
\nabla F=2dg+20(q-q_0),\quad
\nabla^2F=2gg^T+20I_3+2d
\begin{pmatrix}0&0&0\\0&0&-1\\0&-1&0\end{pmatrix}.
\]
从 $q_0$ 出发每步解 Hessian 线性系统, 实际得到
\[
(V,I,R)\approx(2.894121035469170,\ 1.411806724241844,\ 2.008299961656205),
\quad F_{\min}\approx0.005884725643198.
\]
梯度最大绝对值约 $4.5\times10^{-15}$, Hessian 三个特征值约为 $19.9824,20.1123,33.9582$, 所以这是严格局部极小点.

还可验证全局最优性. $F$ 连续且由二次惩罚项控制, 当 $\|q\|\to\infty$ 时趋无穷, 故存在全局最小点. 因 $F(q_0)=0.01$, 每个全局最小点都落在
$|q_i-q_{0,i}|\le\sqrt{0.001}<0.032$ 的盒内. 该盒内
$|d|<0.25$, 故 Hessian 最小特征值至少为 $20-2|d|>19.5$ (秩一项半正定). 盒是凸集, $F$ 在盒内严格凸, 其中至多一个驻点. 实际找到的驻点在盒内, 因而必为唯一全局最小点.
\end{solution}
